% ===================================================================
%  TÁBLA Uimh. 6 — An Teach Istigh. — Troscán. — Bia. — Soilsiú.
%  Traduction 2026 d'apres texto/io, controlee sur texto/fr, texto/en
%  et texto/es. Consigne : voir texto/ga/10-tabla-01.tex.
%
%  Deux notes, toutes deux marquees « (*) », et deux appels : celui de
%  « babo » au bloc c1-12-1, celui de « lambrequino » au bloc c2-01-1.
%
%  LE BLOC c1-06-1 GARDE SON ECART DECLARE : l'ido y porte « (6) » la
%  ou le francais porte « (16) », et c'est le francais qui a raison.
%
%  LE PLUS GROS LEXIQUE DE LA COLONNE. Ici l'irlandais emprunte
%  beaucoup a l'anglais, et il faut l'accepter : « cupard », « prás »,
%  « pram », « tuáille » sont les mots de la maison irlandaise du
%  vingtieme siecle, non des fautes. Forcer un compose celtique la ou
%  la langue emploie l'emprunt serait le meme faux pas que d'inventer
%  du basque pour la marine.
% ===================================================================

%%K t06-apar-1 apar
\VUsaut{6.0mm}
\VUtitre{15.0pt}{60}{TÁBLA Uimh. 6}
\VUsaut{1.4mm}
\VUtitre{11.4pt}{0}{An Teach Istigh, Troscán, Bia,}
\VUsaut{0.4mm}
\VUtitre{11.4pt}{0}{Soilsiú.}
\VUsaut{2.2mm}

%%K t06-apar-2 apar
Tá an teach réidh. Tá uncail Eoin tar éis bogadh isteach ina áit chónaithe nua, a
bhfuil a leagan amach ar eolas againn cheana. Feicimis anois conas a chuir sé
troscán ann. Ar dtús rachaimid go dtí:

%%K t06-tit-1 sub
\VUpk{10.2pt}{40}{An seomra codlata.}
\VUsaut{3.0mm}

%%K t06-c1-01-1 p
1. --- Tháinig muid ag droch-am, mar tá an cailín aimsire ag glanadh an
\VUgras{tseomra}.

%%K t06-c1-02-1 p
2. --- Ar an \VUgras{doirteal níocháin} \textsuperscript{(1)} tá rudaí éagsúla anseo
is ansiúd, a nítear, a chíortar agus a chóirítear leo. Is iad an \VUgras{báisín}
\textsuperscript{(2)} agus an \VUgras{crúsca} \textsuperscript{(3)}, an
\VUgras{scuab ghruaige} \textsuperscript{(4)}, an \VUgras{cíor}
\textsuperscript{(5)}, an \VUgras{gallúnach} \textsuperscript{(6)} agus an
\VUgras{spúinse} \textsuperscript{(7)}, agus ar deireadh an \VUgras{buidéilín}
\textsuperscript{(8)} le huisce béil.

%%K t06-c1-03-1 p
3. --- Ar an urlár, os comhair an doirteail, tá \VUgras{buicéad}
\textsuperscript{(9)} d'uisce salach.

%%K t06-c1-04-1 p
4. --- Sa \VUgras{halla} \textsuperscript{(10)} feictear roinnt \VUgras{crúcaí}
\textsuperscript{(13)} agus dhá \VUgras{scáth fearthainne} \textsuperscript{(11)} i
\VUgras{seastán} \textsuperscript{(12)}.

%%K t06-c1-05-1 p
5. --- Ar an taobh eile den doras tá an \VUgras{cupard scátháin}
\textsuperscript{(14)}, ina gcoinnítear an \VUgras{línéadach}
\textsuperscript{(15)}.

%%K t06-c1-06-1 p
6. --- Bainimis leas as go bhfuil an \VUgras{chailleach seomra}
\textsuperscript{(16)} ag cóiriú na \VUgras{leapa} \textsuperscript{(17)}, agus
scrúdaímis a codanna. Iompraíonn an \VUgras{leaba} \textsuperscript{(20)}
\VUgras{tocht sprionga} maith \textsuperscript{(21)}, ar a gcuirfidh Iustín gan
mhoill an \VUgras{tocht} olla \textsuperscript{(22)}. Ansin tógfaidh sí ó na
cathaoireacha dhá \VUgras{bhraillín} \textsuperscript{(23)} agus an
\VUgras{blaincéad} \textsuperscript{(24)}. Ansin cuirfidh sí ag ceann na leapa an
\VUgras{piliúr fada} \textsuperscript{(25)} agus an \VUgras{piliúr}
\textsuperscript{(26)}, atá in éineacht leis an \VUgras{gcuilt chlúimh}
\textsuperscript{(27)} ar an \VUgras{rug cois leapa} \textsuperscript{(32)}. Leis an
scuab croiticín croithfidh sí na \VUgras{cuirtíní} \textsuperscript{(19)} agus an
\VUgras{ceannbhrat} \textsuperscript{(18)}, agus beidh cuid den obair déanta.

%%K t06-c1-07-1 p
7. --- Ansin beidh uirthi beagán ord a chur ar an \VUgras{mbord oíche}
\textsuperscript{(28)}, an \VUgras{clog aláraim} \textsuperscript{(29)} a thochras,
an \VUgras{solas oíche} \textsuperscript{(30)} a ullmhú agus an \VUgras{choinneal}
\textsuperscript{(31)} a bhreith léi, chun an coinnleoir a ghlanadh.

%%K t06-tit-2 sub
\VUsaut{5.0mm}
\VUpk{10.2pt}{40}{An chiseán oibre agus fearas an teallaigh.}
\VUsaut{3.0mm}

%%K t06-c1-08-1 p
8. --- Tá géarghá le hord ar an \VUgras{gciseán oibre} \textsuperscript{(34)} in
aice leis an \VUgras{stól} \textsuperscript{(33)} freisin. Beidh gá an
\VUgras{bhróidnéireacht} \textsuperscript{(35)} a fhilleadh, an \VUgras{bord
oibre} \textsuperscript{(38)} a líonadh leis an \VUgras{mhéaracán}, leis na
\VUgras{snáitheanna} \textsuperscript{(36)}, leis na \VUgras{snáthaidí}
\textsuperscript{(37)} agus leis an \VUgras{siosúr} \textsuperscript{(39)}, agus
clúdach an bhoird a ghlasáil leis an \VUgras{eochair} \textsuperscript{(40)}.

%%K t06-c1-09-1 p
9. --- Ar an \VUgras{mbord cruinn} \textsuperscript{(41)} i lár an tseomra
athróidh Iustín an t-uisce sa \VUgras{gcarafa} \textsuperscript{(42)}. Cuirfidh sí
an \VUgras{siúcra} sa \VUgras{siúcrán} \textsuperscript{(43)}, glanfaidh sí na
\VUgras{maidí siúcra} \textsuperscript{(44)} agus béarfaidh sí léi an \VUgras{scuab
éadaigh} \textsuperscript{(46)}.

%%K t06-c1-10-1 p
10. --- Tabharfaidh taobh dheis an tseomra níos lú oibre di, mar tá an
\VUgras{tolg} \textsuperscript{(47)} agus a \VUgras{chúisín}
\textsuperscript{(48)} ina n-áit, agus ó tharla nach bhfuil sé fuar, níor baineadh
le fearas an \VUgras{teallaigh} \textsuperscript{(49)} ach oiread. Tá an
\VUgras{sluaistín} \textsuperscript{(50)}, an \VUgras{teanchair}
\textsuperscript{(51)}, na \VUgras{cinn tine} \textsuperscript{(55)} agus an
\VUgras{ghráta luaithe} \textsuperscript{(56)} glan; níl an \VUgras{scáileán
teallaigh} \textsuperscript{(54)} bogtha, agus tá an \VUgras{bosca adhmaid}
\textsuperscript{(52)} lán de \VUgras{lomáin} \textsuperscript{(53)}.

%%K t06-noto-1 noto
\VUnotes{143.2mm}{%
(*) Babo = leanbh beag; Béarla \textit{baby}, Fr. \textit{bébé}, Iod.
\textit{bambino}.%
}

%%K t06-noto-2 noto
\VUnotes{143.2mm}{%
(*) Lambrequino = Fr. \textit{lambrequin}, Sp. \textit{lambrequines}, Port.
\textit{lambrequins}, agus fiú Gearm. agus Béarla \textit{lambrequins} --- is é
sin le rá i nGearmáinis, i mBéarla, i bhFraincis, i bPortaingéilis agus i
Spáinnis.%
}

%%K t06-c1-11-1 p
11. --- Beidh uirthi fós an \VUgras{clog teallaigh} \textsuperscript{(57)}, na
\VUgras{coinnleoirí} \textsuperscript{(58)} agus na \VUgras{tráidirí do rudaí
beaga} \textsuperscript{(59)} a chroitheadh, agus ansin an \VUgras{scáthán}
\textsuperscript{(60)} a chuimilt.

%%K t06-c1-12-1 p
12. --- Maidir le \VUgras{cliabhán} \textsuperscript{(61)} an linbh (babo (*)), tá
sé réidh cheana.

%%K t06-tit-3 sub
\VUsaut{5.0mm}
\VUpk{10.2pt}{40}{An seomra suí.}
\VUsaut{3.0mm}

%%K t06-c2-01-1 p
1. --- Seo an seomra suí. Is \VUgras{seomra} an-álainn é. Tá an teallach sa chúinne
deas maisithe le \VUgras{dealbhóg} ghalánta \textsuperscript{(1)} agus le dhá
\VUgras{vása} áille \textsuperscript{(3)}, agus é clúdaithe le \VUgras{lambraicín}
veilbhite \textsuperscript{(2)} (*).

%%K t06-c2-02-1 p
2. --- Chun nach lasfadh na spréacha ón tine an rug, cuireadh os comhair an
teallaigh \VUgras{scáileán} práis \textsuperscript{(4)}.

%%K t06-c2-03-1 p
3. --- Lena thaobh tá \VUgras{deasc} ghalánta bhan ar oscailt \textsuperscript{(5)}.
Anseo a scríobhann \VUgras{bean an tí} \textsuperscript{(23)} a cuid litreacha.

%%K t06-c2-04-1 p
4. --- Tá an fhuinneog maisithe le \VUgras{cuirtíní tulla} \textsuperscript{(6)},
ceangailte le \VUgras{ceangail} \textsuperscript{(8)} ar \VUgras{crúcaí}
\textsuperscript{(7)}, agus le cuirtíní troma éadaigh le \VUgras{lambraicín}
\textsuperscript{(9)} thuas.

%%K t06-c2-05-1 p
5. --- I gcuas na fuinneoige coinníonn \VUgras{seastán bláthanna}
\textsuperscript{(11)} \VUgras{planda} glas \textsuperscript{(10)}, pailm iontach a
shroicheann a duilleoga beagnach an \VUgras{lampa mine}
\textsuperscript{(12)}.

%%K t06-c2-06-1 p
6. --- Is í Máire a rinne an \VUgras{scáthlán} \textsuperscript{(13)}, an
\VUgras{iníon uasal} ghleoite \textsuperscript{(14)} ag an \VUgras{bpianó}
\textsuperscript{(15)}. Ina suí go díreach ar an \VUgras{stól}
\textsuperscript{(16)}, tá sí ag foghlaim píosa ceoil. Tá sí an-oilte agus
an-néata, mar a thaispeánann a \VUgras{cófra ceoil} \textsuperscript{(17)}: tá ord
iomlán ann.

%%K t06-tit-4 sub
\VUsaut{5.0mm}
\VUpk{10.2pt}{40}{An Dailtín.}
\VUsaut{3.0mm}

%%K t06-c2-07-1 p
7. --- Tá a deartháir Pól ag lorg \VUgras{leabhair} \textsuperscript{(19)} sa
\VUgras{leabhragán} \textsuperscript{(18)}. Is \VUgras{buachaill}
\textsuperscript{(20)} corrthónach agus do-fhulaingthe ar fad é. Ar crochadh den
chófra chun leabhar atá ró-ard a bhaint, tá contúirt ann go leagfaidh sé an
\VUgras{busta} \textsuperscript{(21)} atá thuas.

%%K t06-c2-08-1 p
8. --- Baineann sé leas don amaidí seo as go bhfuil an mháthair gafa i gcomhrá le
cara, \VUgras{bean uasal} \textsuperscript{(22)} a tháinig chun fanacht cúpla lá
aici. Tá an mháthair ina suí ar an \VUgras{tolg} \textsuperscript{(24)} in aice
leis an \VUgras{gcathaoir uilleann} \textsuperscript{(25)}, agus a cara ar
\VUgras{chathaoir} \textsuperscript{(27)}, nach bhfeicimid di ach an
\VUgras{droim} \textsuperscript{(26)}.

%%K t06-c2-09-1 p
9. --- Sa \VUgras{fhráma} mór \textsuperscript{(30)} ar an mballa tá
\VUgras{portráid} \textsuperscript{(29)} gaoil. San \VUgras{albam}
\textsuperscript{(32)}, atá ar an mbord, tá \VUgras{grianghraif}
\textsuperscript{(33)} an chuid eile de na gaolta bailithe. Tá na páistí díreach tar
éis é a chasadh, mar tá na \VUgras{cathaoireacha} \textsuperscript{(31)} fós
bailithe timpeall an bhoird.

%%K t06-c2-10-1 p
10. --- Tugadh an \VUgras{vása Síneach} poircealláin \textsuperscript{(34)} in aice
leis an \VUgras{scian pháipéir} \textsuperscript{(35)} do Mháire ar a lá ainm, agus
is é a huncail a thug ón tSín an \VUgras{scáileán} \textsuperscript{(36)} a
líonann cúinne an tseomra.

%%K t06-c2-11-1 p
11. --- Beoite ag na \VUgras{pictiúir} áille \textsuperscript{(37)} ar an
\VUgras{mballa} \textsuperscript{(42)}, soilsithe ag an \VUgras{gcoinnleoir craobhach}
síleála \textsuperscript{(38)}, a soilsíonn a \VUgras{bhoilgeoga leictreacha}
\textsuperscript{(39)} chomh meidhreach sin sa tráthnóna, dealraíonn an seomra suí
seo an-taitneamhach. Déanann an \VUgras{rug} bog \textsuperscript{(40)}, atá
leata ar an urlár, agus an \VUgras{cloigín leictreach} chomh háisiúil sin
\textsuperscript{(41)} cluthar ar fad é.

%%K t06-tit-5 sub
\VUsaut{3.6mm}
\VUpk{10.2pt}{40}{An seomra bia.}
\VUsaut{3.0mm}

%%K t06-c3-01-1 p
1. --- Buaileadh an clog don dinnéar. Tá an teaghlach ar fad, seachas Máire,
bailithe timpeall an bhoird. Tá an seomra bia téite go deas ag \VUgras{sorn}
tíleanna den scoth \textsuperscript{(1)} lena \VUgras{phíopa} caol
\textsuperscript{(2)}.

%%K t06-c3-02-1 p
2. --- Níl mórán troscáin sa seomra seo. Feadh an bhalla, os cionn an
\VUgras{stóil} \textsuperscript{(4)}, tá \VUgras{tobar níocháin} maisiúil
\textsuperscript{(3)} ar crochadh. Lena thaobh tá an \VUgras{cófra bia}
\textsuperscript{(5)}, ar a gcuirtear na miasa fuara, an mhilseog agus na
gréithe éagsúla nach bhfuil ag teastáil fós.

%%K t06-c3-03-1 p
3. --- Mar sin, ar an tseilf uachtarach feicimid \VUgras{sicín rósta}
\textsuperscript{(7)}, \VUgras{iasc} \textsuperscript{(8)} agus \VUgras{vása}
\textsuperscript{(10)} le \VUgras{torthaí} \textsuperscript{(9)}. Níos ísle tá an
\VUgras{cháis} \textsuperscript{(11)}, an \VUgras{t-arán} \textsuperscript{(12)}
agus an \VUgras{gléas boird} \textsuperscript{(13)}, ina bhfuil gach a bhfuil de
dhíth chun an sailéad a bhlaisiú: ola, fínéagar, \VUgras{salann}
\textsuperscript{(14)} agus \VUgras{piobar} \textsuperscript{(15)}. Ar deireadh, ar
an tseilf íochtarach tá \VUgras{pióg} \textsuperscript{(17)} in aice leis an
\VUgras{mustardán} \textsuperscript{(16)}.

%%K t06-c3-04-1 p
4. --- Tá clog an tseomra bia, ar a dtugtar an \VUgras{clog balla}
\textsuperscript{(20)}, de chruth an-neamhghnách. Tá sé curtha i bhfráma adhmaid,
ina bhfuil freisin an \VUgras{baraiméadar} \textsuperscript{(19)} agus an
\VUgras{teirmiméadar} \textsuperscript{(18)}.

%%K t06-tit-6 sub
\VUsaut{4.6mm}
\VUpk{10.2pt}{40}{Conas a thugtar an dinnéar.}
\VUsaut{3.0mm}

%%K t06-c3-05-1 p
5. --- Timpeall an tseomra ar fad tá na ballaí clúdaithe le \VUgras{painéil
adhmaid} \textsuperscript{(21)} de dhair chéarach. Dúnann \VUgras{cuirtín}
éadaigh \textsuperscript{(22)} go maith go leor an doras a théann go dtí an
vearanda. Ann sin a rachaidh an teaghlach tar éis an dinnéir chun caife a ól. Tá na
\VUgras{cupáin} \textsuperscript{(24)} agus na \VUgras{sásair}
\textsuperscript{(25)} leagtha amach cheana ar an \VUgras{drisiúr}
\textsuperscript{(23)}.

%%K t06-c3-06-1 p
6. --- Os cionn an bhoird tá \VUgras{lampa crochta} \textsuperscript{(26)}
ceangailte den tsíleáil; líonann a solas an-gheal an seomra bia ar fad sa
tráthnóna.

%%K t06-c3-07-1 p
7. --- Féachaimis anois ar an \VUgras{mbord dinnéir} féin \textsuperscript{(27)}.
Nuair a tháinig an teaghlach isteach, bhí an bord leagtha. Ar an \VUgras{éadach
boird} \textsuperscript{(28)} cuireadh ag gach áit \VUgras{pláta}
\textsuperscript{(33)}, \VUgras{naipcín} \textsuperscript{(29)},
\VUgras{spúnóg} \textsuperscript{(34)}, \VUgras{forc} \textsuperscript{(35)},
\VUgras{scian} \textsuperscript{(36)} agus \VUgras{gloine}
\textsuperscript{(32)}. Bhí ann freisin \VUgras{buidéil} \textsuperscript{(30)} don
fhíon agus \VUgras{carafa} \textsuperscript{(31)} don uisce.

%%K t06-c3-08-1 p
8. --- Thug an \VUgras{seirbhíseach} \textsuperscript{(39)} isteach ar dtús an
\VUgras{sáspan anraith} \textsuperscript{(37)}, ina bhfuil an t-anraith fós ag
gail. Anois tá sé ag tabhairt an chéad \VUgras{mhias} \textsuperscript{(38)}, an
ceann feola. Ansin freastalóidh sé ar chách leis na cinn atá ainmnithe againn; agus
ar deireadh, tar éis na \VUgras{milseoige}, bainfidh sé leas as go bhfuil na
tiarnaí tar éis dul chuig an vearanda chun an bord a ghlanadh agus ord a chur ar an
seomra bia arís.

%%K t06-c3-08-2 p
\textit{Nóta.} --- \textit{Níl na focail choitianta faoin mbia tugtha anseo ach i
línte ginearálta. Cuirfear leis an liosta sa dá thábla «An tOstán» (uimh. 13) agus
«An Margadh» (uimh. 15).}

%%K t06-tit-7 sub
\VUsaut{4.6mm}
\VUpk{10.2pt}{40}{An chistin.}
\VUsaut{3.0mm}

%%K t06-c4-01-1 p
1. --- Sa chistin, mar a bhféachaimid tríd an doras leathoscailte, buaileann
mí-ord pictiúrtha linn. Os comhair an \VUgras{teallaigh} \textsuperscript{(1)}, mar
a bhfuil an \VUgras{tine} \textsuperscript{(3)} ag cnagadh, tá an \VUgras{meaisín
bhior} \textsuperscript{(5)} suite. Tá \VUgras{rósta} iontach \textsuperscript{(7)}
curtha ar an \VUgras{mbior} \textsuperscript{(6)} agus é ag bruith.

%%K t06-c4-02-1 p
2. --- D'fhan an \VUgras{boilg} \textsuperscript{(8)} agus an \VUgras{sluaistín
guail} \textsuperscript{(9)}, ar baineadh úsáid astu ar ball, ar an urlár in
éineacht leis na \VUgras{lomáin} \textsuperscript{(10)}.

%%K t06-c4-03-1 p
3. --- Ar sheilf an teallaigh tá ord: an \VUgras{lóchrann}
\textsuperscript{(11)}, an \VUgras{bosca cipíní} \textsuperscript{(13)} lena
\VUgras{chipíní} \textsuperscript{(12)}, an \VUgras{coinnleoir}
\textsuperscript{(14)} agus an \VUgras{coinnleoir oíche} \textsuperscript{(15)}
lena \VUgras{choinneal} \textsuperscript{(16)} --- gach rud ina áit féin. Os a
choinne sin, fágadh an \VUgras{buicéad guail} mór \textsuperscript{(18)}, lán de
\VUgras{ghual} \textsuperscript{(17)}, ar a ndeachaigh an \VUgras{cócaire}
\textsuperscript{(23)} síos sa siléar ar ball, i lár na cistine.

%%K t06-c4-04-1 p
4. --- Is é an fáth go gcaithfidh an bhean bhocht deifir a dhéanamh, chun go mbeidh
an bricfeasta in am. Anois tá sí ag an \VUgras{sorn} \textsuperscript{(19)}, ag
suaitheadh an \VUgras{anlainn} \textsuperscript{(24)}, nach féidir a ligean dó
dó.

%%K t06-c4-05-1 p
5. --- San \VUgras{oigheann} \textsuperscript{(20)} tá an phióg a rinne sí do na
páistí don sólaistí á bácáil. Os cionn an tsoirn tá an \VUgras{ladar}
\textsuperscript{(21)} agus an \VUgras{spúnóg pholltach} \textsuperscript{(22)} ar
crochadh.

%%K t06-tit-8 sub
\VUsaut{4.0mm}
\VUpk{10.2pt}{40}{Gréithe na cistine.}
\VUsaut{3.0mm}

%%K t06-c4-06-1 p
6. --- Tá an \VUgras{tsraith corcán} \textsuperscript{(25)}, atá ar crochadh ar chúl
an tseomra, an-ghlan. Tá na \VUgras{corcáin} phráis áille \textsuperscript{(26)}, an
\VUgras{crúsca bainne} \textsuperscript{(27)}, an \VUgras{síothlán}
\textsuperscript{(28)} agus an \VUgras{scríobán} \textsuperscript{(29)} glanta go
cúramach.

%%K t06-c4-07-1 p
7. --- Tá an \VUgras{sáiléar} \textsuperscript{(30)} ar sheilfín bheag in aice leis
an \VUgras{taephota} \textsuperscript{(31)} agus leis an \VUgras{mbuidéal ola}
\textsuperscript{(32)}. Sa \VUgras{tarraiceán} \textsuperscript{(33)}, is dócha, tá
na gréithe boird agus sceana na cistine.

%%K t06-c4-08-1 p
8. --- Tá an \VUgras{scuab} \textsuperscript{(34)} agus an \VUgras{scuab chleití}
\textsuperscript{(35)} sa chúinne. Tá an cócaire díreach tar éis úsáid a bhaint as
an \VUgras{gceap gearrtha} \textsuperscript{(36)}; d'fhág sí in aice leis an
bhfuinneog é.

%%K t06-c4-09-1 p
9. --- Tá an \VUgras{tuáille} \textsuperscript{(37)} ar crochadh ar an mballa in
aice leis an \VUgras{tonnadóir} \textsuperscript{(38)}. Sa chuid seo den chistin
coinnítear an \VUgras{greille} \textsuperscript{(39)}, an \VUgras{tua bheag}
\textsuperscript{(40)} agus an \VUgras{clár gearrtha} \textsuperscript{(41)}.
Ansin, os cionn na tua bige, ar an \VUgras{tseilf} \textsuperscript{(42)}, tá na
\VUgras{pota} \textsuperscript{(43)}, an \VUgras{corcán iarainn}
\textsuperscript{(44)} lena \VUgras{chlúdach} \textsuperscript{(45)} agus an
\VUgras{coire} \textsuperscript{(46)}. Tá an \VUgras{friochtán}
\textsuperscript{(47)} ar crochadh lena thaobh.

%%K t06-c4-10-1 p
10. --- Níl ach an \VUgras{sconna} práis \textsuperscript{(48)} ag caitheamh
loinnre sa chúinne seo. Caithfidh an \VUgras{doirteal} \textsuperscript{(49)}, ar a
bhfuil \VUgras{crúsca} mór \textsuperscript{(50)}, bheith an-áisiúil lena chófra
beag, mar a gcoinnítear an \VUgras{bairillín fínéagair} \textsuperscript{(51)} lena
\VUgras{sconna beag} \textsuperscript{(52)}, an \VUgras{bosca bruscair}
\textsuperscript{(53)} agus na \VUgras{gréithe salacha} \textsuperscript{(54)}.

%%K t06-tit-9 sub
\VUsaut{4.0mm}
\VUpk{10.2pt}{40}{Cad eile a choinnítear sa chistin.}
\VUsaut{3.0mm}

%%K t06-c4-11-1 p
11. --- Is dócha go mbaineann an \VUgras{tobán} mór \textsuperscript{(55)}, ina
nitear na \VUgras{glasraí} \textsuperscript{(58)}, agus an \VUgras{coire} iarainn
\textsuperscript{(56)}, atá ar an \VUgras{urlár tíleanna} \textsuperscript{(57)},
leis an gcófra beag sin freisin.

%%K t06-c4-12-1 p
12. --- Ní easpa sorn atá ar an gcócaire, mar seo, ar chúinne an bhoird, tá
\VUgras{dóiteoir gáis} \textsuperscript{(59)}, agus air tá uisce á théamh i
gcorcán. Insíonn an \VUgras{gal} \textsuperscript{(60)} a thagann as dúinn go
gcaithfidh an t-uisce bheith ag fiuchadh. Is dócha gur don chaife an t-uisce sin,
mar seo ar an gceann eile den bhord tá an \VUgras{pota caife}
\textsuperscript{(64)} agus an \VUgras{muileann caife} \textsuperscript{(63)}.

%%K t06-c4-13-1 p
13. --- Tá an dóiteoir ceangailte den \VUgras{phíopa gáis} \textsuperscript{(62)} le
\VUgras{feadán rubair} fada \textsuperscript{(61)}.

%%K t06-c4-14-1 p
14. --- Tá an \VUgras{bhean lae} \textsuperscript{(66)}, a thagann chun cabhrú leis
an gcócaire, ag ullmhú na nglasraí don dinnéar. Nuair a bheidh siad glanta agus
nite, cuirfidh sí sa \VUgras{gcorcán} \textsuperscript{(65)} iad go dtí go dtiocfaidh
an t-am iad a bhruith.

%%K t06-tit-10 sub
\VUsaut{4.0mm}
{\centering\fontsize{10.2pt}{10.2pt}\selectfont\textls[40]{\scshape An seomra folctha.} \textit{(An seomra ina bhfolctar.)}\par}
\VUsaut{3.0mm}

%%K t06-c5-01-1 p
1. --- Níl fágtha ach áit amháin le feiceáil, chun príomhsheomraí an tí a aithint:
leagan amach an tseomra folctha. Ní bheidh sé fada, mar tá sé beag, agus feicfimid
go tapa go bhfuil \VUgras{téitheoir} maith ag an \VUgras{bhfolcadán}
\textsuperscript{(1)} \textsuperscript{(2)}, go bhfuil an \VUgras{gallúnadán}
\textsuperscript{(3)} ar láimh ag an duine a fholcann, agus go gcaithfidh an
\VUgras{crochadán tuáillí} \textsuperscript{(5)} triomú na \VUgras{dtuáillí}
\textsuperscript{(4)} a éascú.

%%K t06-c5-02-1 p
2. --- Tá ballaí an tseomra seo clúdaithe le \VUgras{tíleanna gloinithe}
\textsuperscript{(6)}, rud a lig dóibh an \VUgras{cith} \textsuperscript{(7)} a
chur, gan eagla go millfeadh an t-uisce a spraeáiltear iad. Críochnaíonn
\VUgras{báisín} beag since \textsuperscript{(8)}, do fholcadh na gcos, an
troscán.
\VUsaut{6.0mm}
\VUfilet{22mm}
